\chapter{A Rust futtatókörnyezet, útvonalkeresés és vizualizáció}

\section{A Rust gráfállapot}

A Rust réteg a gráfot közvetlenül a meccs-JSON állományokból és a játékos-adatbázisból építi fel. A \textit{GraphState} a szomszédsági listák mellett tárolja a \textit{pair\_relations} szerkezetet, a játékosprofilok gyors elérését, a klasztertagságot, a landmark-távolságokat, a minimumköltségeket és az aktív datasethez tartozó snapshot-metaadatot.

\section{Útvonalkereső algoritmusok}

A jelenlegi Rust futtatókörnyezet négy algoritmust támogat: BFS, Dijkstra, kétirányú keresés és A*. Az algoritmusok két élértelmezésben futhatnak: \textit{social-path} csak szövetségesi kapcsolatokon, \textit{battle-path} összes ismétlődő kapcsolaton.

Az élköltség-modell: súlyozatlan módban minden él költsége fix $1000$. Súlyozott módban:
\begin{equation}
c(u,v)=\left\lceil \frac{1000}{w(u,v)} \right\rceil
\end{equation}
ahol $w(u,v)$ az adott útvonalmódnak megfelelő kapcsolaterősség.

\begin{table}[H]
\centering
\caption{A Pathfinder Labban szereplő algoritmusok összevetése}
\begin{tabularx}{\textwidth}{>{\raggedright\arraybackslash}p{2.8cm} >{\raggedright\arraybackslash}p{2.8cm} X}
\toprule
\textbf{Algoritmus} & \textbf{Mikor természetes?} & \textbf{Megjegyzés a projektben} \\
\midrule
BFS & súlyozatlan legrövidebb út & Jó baseline, könnyen magyarázható \\
Dijkstra & súlyozott utak & Referenciaalgoritmus a költség-alapú módokban \\
Kétirányú keresés & nagyobb gráf, kevésbé irányított keresés & Gyakorlati gyorsítás bizonyos esetekben \\
A* & jól kialakított alsó becsléssel & A projektben egzakt marad, nem csak vizuális heurisztika \\
\bottomrule
\end{tabularx}
\end{table}

\section{Egzakt A* és ALT heurisztika}

Az ALT (A* + Landmarks + Triangle inequality) heurisztika \cite{goldberg2005} elvéből indul ki: ha $L$ egy landmark-csomópont, akkor
\begin{equation}
d(u, t) \geq |d(L, t) - d(L, u)|
\end{equation}
minden $L$ landmarkra. A landmark-távolságok a \textit{GraphState} felépítésekor előre kiszámítódnak minden \textit{ModeKey} kombinációra, és az összes keresési kérésre újrafelhasználhatók. Az A* helyes legrövidebb utat ad; egységtesztek igazolják, hogy az A* eredménye egyezik a Dijkstra-eredménnyel minden releváns esetben.

\section{Pathfinder Lab}

A Pathfinder Lab az a réteg, ahol a Rust algoritmusok, a backend-szerződések és a böngészőben értelmezhető interakció egyetlen demonstrációvá állnak össze.

\begin{figure}[H]
\centering
\includegraphics[width=0.98\textwidth,height=0.74\textheight,keepaspectratio]{pathfinder-replay-panel.png}
\caption{A Pathfinder Lab replay panele -- az útvonalösszegző, a léptetési vezérlők és a meglátogatott csomópontok statisztikái}
\end{figure}

Az alábbi mérési eredmény különösen szemléletes:

\begin{table}[H]
\centering
\caption{BFS és súlyozott keresés összehasonlítása -- ugyanazon játékospár, eltérő közelségi fogalmak}
\begin{tabularx}{\textwidth}{>{\raggedright\arraybackslash}p{3cm} >{\raggedright\arraybackslash}p{3cm} >{\raggedright\arraybackslash}p{2cm} X}
\toprule
\textbf{Algoritmus} & \textbf{Metrika} & \textbf{Eredmény} & \textbf{Értelmezés} \\
\midrule
BFS & Minimális ugrásszám & 3 & Legrövidebb szociális lánc \\
Dijkstra & Min. súlyozott költség & 8 & Legerősebb szövetségesi útvonal \\
A* & Min. súlyozott költség & 8 & Egyezik Dijkstrával \\
Kétirányú & Min. ugrásszám & 3 & Egyezik BFS-sel \\
\bottomrule
\end{tabularx}
\end{table}

A BFS által talált 3-ugrás-os útvonal strukturálisan létezik, de gyenge vagy egyszeri meccsből eredő kapcsolatokon halad át. A Dijkstra súlyozott módban elkerüli ezeket, és helyette egy 8-ugrás-os útvonalat talál, amelyen minden él ismétlődő szövetségesi viszony.

\begin{table}[H]
\centering
\caption{A négy algoritmus összehasonlítása a legfontosabb dimenziók mentén}
\begin{tabularx}{\textwidth}{>{\raggedright\arraybackslash}p{3cm} >{\raggedright\arraybackslash}p{2.5cm} >{\raggedright\arraybackslash}X >{\raggedright\arraybackslash}X}
\toprule
\textbf{Algoritmus} & \textbf{Optimális?} & \textbf{Memória} & \textbf{Fő erősség a projektben} \\
\midrule
BFS & Igen (súlyozatlan) & $O(n)$ & Egyszerű baseline, jól magyarázható replay \\
Dijkstra & Igen (súlyozott) & $O(n + m)$ & Referencia-algoritmus, lexikografikus döntetlen-kezelés \\
Kétirányú & Igen (súlyozatlan) & $O(n)$ & Potenciális gyorsulás klaszteres gráfon \\
A* & Igen (heurisztika elfogadható) & $O(n)$ & Célirányos bejárás, egzakt eredmény \\
\bottomrule
\end{tabularx}
\end{table}

\section{A 3D globális vizualizáció}

A \textit{bird's-eye sphere} a teljes hálózat vizuális áttekintésére szolgál. A Rust réteg előre exportál minden fontos artifaktumot, és a frontend csak renderel.

\begin{table}[H]
\centering
\caption{A bird's-eye export fő artifaktumai}
\begin{tabularx}{\textwidth}{>{\raggedright\arraybackslash}p{4.2cm}X}
\toprule
\textbf{Fájl} & \textbf{Szerep} \\
\midrule
\textit{manifest.json} & Metaadatok, darabszámok, skálázási és generálási információk \\
\textit{node\_meta.json} & Csomópontszintű leíró metaadatok, frontend-hozzárendelések \\
\textit{node\_positions.f32} & Lebegőpontos koordináták tömbje a gyors rendereléshez \\
\textit{node\_metrics.u32} & Tömörített numerikus csomóponti mutatók \\
\textit{edge\_pairs.u32} & Él-végpontpárok tömbje \\
\textit{edge\_props.u32} & Élhez tartozó kódolt tulajdonságok \\
\bottomrule
\end{tabularx}
\end{table}

\begin{figure}[H]
\centering
\includegraphics[width=0.95\textwidth,keepaspectratio]{graph artifact package dataflow.pdf}
\caption{A Graph V2 artifaktumcsomag adatfolyama: a Rust-exportőr bináris és JSON-rétegeket ír, az Express API kiszolgálja, a frontend WebGL-renderelőbe fogyasztja}
\end{figure}

\begin{figure}[H]
\centering
\includegraphics[width=\textwidth,height=0.82\textheight,keepaspectratio]{graph artifact components.pdf}
\caption{A Graph V2 pipeline bináris és JSON artifact-fájljainak szerkezete}
\end{figure}

\begin{figure}[H]
\centering
\includegraphics[width=0.98\textwidth,height=0.72\textheight,keepaspectratio]{birds-eye-current-state.png}
\caption{A teljes hálózat 3D madártávlati nézete az előre generált Rust-artifaktumokra építve}
\end{figure}

\begin{figure}[H]
\centering
\includegraphics[width=0.98\textwidth,height=0.72\textheight,keepaspectratio]{flexset_associative_graph_clusters.png}
\caption{Sűrűbb klaszterfoltok és a térbeli elkülönítés a globális nézetben}
\end{figure}

\section{A strukturális egyensúly mint módszertani határeset}

A repositoryban megvalósított \textit{signed\_balance.rs} modul determinisztikusan előállít egy előjeles gráfot és triádokat enumerál. A visszaminősítés oka szemantikai: az \textit{enemyWeight} dominanciája technikailag negatív élként jelölhető, de ez nem azonos ellenségességgel. Az ismételt szembekerülés hasonló MMR-sávból, azonos napszakból, a Master$+$ Flex Queue kis aktív játékospoolából, matchmaking-ismétlésből vagy az adatgyűjtési stratégia torzításából is adódhat. A modul kísérleti opcióként megmaradhat; a dolgozat fő pillérei az asszociatív gráfértelmezés, az asszortativitási eredmények és a Brandes-betweenness centralitás.
